\documentclass[a4paper,11pt]{article}
        \usepackage[top=15mm,bottom=18mm,left=16mm,right=16mm]{geometry}
        \usepackage{fontspec}
        \usepackage{xeCJK}
        \setmainfont{Times New Roman}
        \setCJKmainfont{Yu Gothic}
        \setmonofont{Consolas}
        \setCJKmonofont{Yu Gothic}
        \usepackage{booktabs}
        \usepackage{longtable}
        \usepackage{tabularx}
        \usepackage{array}
        \usepackage{enumitem}
        \usepackage{hyperref}
        \usepackage{listings}
        \usepackage{xcolor}
        \hypersetup{
          colorlinks=true,
          linkcolor=blue,
          urlcolor=blue,
          pdftitle={自動校正ロジック関数群},
          pdfauthor={Codex}
        }
        \lstset{
          basicstyle=\ttfamily\small,
          breaklines=true,
          columns=fullflexible,
          keepspaces=true,
          frame=single,
          framerule=0.2pt,
          rulecolor=\color{black},
          showstringspaces=false
        }
        \setlength{\parskip}{0.42em}
        \setlength{\parindent}{1em}
        \renewcommand{\arraystretch}{1.15}
        \title{25. 自動校正ロジック関数群}
        \author{solar\_ws0129-main}
        \date{2026-07-29}
        \begin{document}
        \maketitle

        \section*{対象}
        \begin{tabularx}{\linewidth}{>{\raggedright\arraybackslash}p{0.18\linewidth} X}
        \toprule
        項目 & 内容 \\
        \midrule
        ファイル & \path|mpc_solarcar/solar_autocal_logic.py| \\
        種別 & Python \\
        区分 & runtime helper \\
        役割 & solar\_autocal\_node が使う昼間停止時 aux 推定などの純ロジック関数をまとめる。 \\
        起動文脈 & Node から切り出された小さな判定ロジック。 \\
        \bottomrule
        \end{tabularx}

        \section*{このファイルを読む理由}
        solar\_autocal\_node が使う昼間停止時 aux 推定などの純ロジック関数をまとめる。

        \begin{itemize}[leftmargin=1.6em]
        \item Node 依存を持たないので単体試験しやすい。
        \end{itemize}

        \section*{呼び出し関係}
        \subsection*{どこから呼ばれるか}
        \begin{itemize}[leftmargin=1.6em]
        \item \path|mpc_solarcar/solar_autocal_node.py|
        \end{itemize}

        \subsection*{次に読むと理解が進むファイル}
        \begin{itemize}[leftmargin=1.6em]
        \item 特記事項なし。
        \end{itemize}

        \subsection*{同一リポジトリ内の主要依存}
        \begin{itemize}[leftmargin=1.6em]
        \item 同一リポジトリ内の明示 import は少ない。
        \end{itemize}

        \section*{ソース冒頭の把握}
        モジュール docstring または先頭意図の要約:

        Pure helpers for bounded online solar-car calibration.

        \subsection*{代表コード断片}
        \begin{lstlisting}[language=Python]
import math


def daytime_stationary_aux_estimate(
    *,
    ghi_wm2: float,
    day_ghi_threshold_wm2: float,
    speed_kmh: float,
    stationary_speed_kmh: float,
    pack_power_w: float,
    solar_power_w: float,
) -> float | None:
    values = (ghi_wm2, speed_kmh, pack_power_w, solar_power_w)
    if not all(math.isfinite(float(value)) for value in values):
        return None
    if float(ghi_wm2) < float(day_ghi_threshold_wm2):
        return None
    if abs(float(speed_kmh)) > float(stationary_speed_kmh):
        return None
    return max(0.0, float(pack_power_w) + float(solar_power_w))
\end{lstlisting}


        \section*{主要構造の読み取り}
        主要関数は daytime\_stationary\_aux\_estimate。

        \subsection*{主要クラス・関数}
\begin{longtable}{p{0.07\linewidth} p{0.16\linewidth} p{0.25\linewidth} p{0.40\linewidth}}
\toprule
行 & 種別 & 名前 & 補足 \\
\midrule
8 & function & \path|daytime_stationary_aux_estimate| &  \\
\bottomrule
\end{longtable}

        \subsection*{ファイルを上から読んだときの定義順}
\begin{longtable}{p{0.07\linewidth} p{0.84\linewidth}}
\toprule
行 & 内容 \\
\midrule
8 & 関数 daytime\_stationary\_aux\_estimate を定義する。 \\
\bottomrule
\end{longtable}

        \subsection*{import 群の詳細}
            \begin{longtable}{p{0.07\linewidth} p{0.33\linewidth} p{0.50\linewidth}}
            \toprule
            行 & import 文 & このファイルで読む理由 \\
            \midrule
            3 & \path|from __future__ import annotations| & 型ヒントの遅延評価を有効にし、自己参照型や forward reference を安全に扱うため。 このファイル内での使用位置は少ないか、間接利用である。 \\
5 & \path|import math| & clip、sqrt、sin/cos、有限判定など数値ロジックに使うため。 このファイル内での主な使用位置は L18。 \\
            \bottomrule
            \end{longtable}

        \section*{関数・クラスを上から順に解説}
        \subsection*{L8 関数 \path|daytime_stationary_aux_estimate|}
            \begin{itemize}[leftmargin=1.6em]
              \item 定義: \path|daytime_stationary_aux_estimate(ghi_wm2, day_ghi_threshold_wm2, speed_kmh, stationary_speed_kmh, pack_power_w, solar_power_w)|
              \item docstring: 明示されていない。
              \item このブロックが直接呼ぶ主な関数・メソッド: abs, all, float, isfinite, max
              \item 戻り値の要点: max(0.0, float(pack\_power\_w) + float(solar\_power\_w)) / None / None / None
            \end{itemize}
            \begin{enumerate}[leftmargin=1.8em]
            \item values に (ghi\_wm2, speed\_kmh, pack\_power\_w, solar\_power\_w) の結果を代入する。
\item 条件 not all((math.isfinite(float(value)) for value in values)) を判定し、真なら内部処理を行う。
\item 条件 float(ghi\_wm2) < float(day\_ghi\_threshold\_wm2) を判定し、真なら内部処理を行う。
\item 条件 abs(float(speed\_kmh)) > float(stationary\_speed\_kmh) を判定し、真なら内部処理を行う。
\item max(0.0, float(pack\_power\_w) + float(solar\_power\_w)) を返す。
            \end{enumerate}











        \section*{処理の流れ}
        \begin{enumerate}[leftmargin=1.7em]
        \item ソース中の主要関数を通じて処理を進める。
        \end{enumerate}

        \section*{このファイルの位置づけ}
        この資料群では、\path|mpc_solarcar/solar_autocal_logic.py| を
        \textbf{runtime helper}の中核として扱う。
        したがって、ここで扱う処理は単独で完結せず、
        前段の profile / launch / telemetry と後段の planner / logger / report へ繋がる。

        \end{document}
